\cotefile{136}
\scanpage{142}
\mspage{2}
\begin{center}- 134 -\end{center}
\margin{IX}

et, en vertu de la structure connue de $I$ ($I\qquad$), une suite exacte
\[
(11.3.3)\qquad 1\longrightarrow R\longrightarrow I\longrightarrow T\longrightarrow 1\quad,
\]
avec $R$ un groupe pro-fini \uline{premier à $\ell$}. Si maintenant nous \uline{partons} de données (11.3.2) et (11.3.3), et introduisons les groupes
\[
(11.3.4)\qquad \Phi=(U\otimes D_\ell)^I/U^I\otimes D_\ell\ ,\quad
\Phi'=(U'\otimes D_\ell)^I/U'^I\otimes D_\ell\ ,\quad
(D_\ell=\mathbb{Q}_\ell/\mathbb{Z}_\ell)
\]
nous allons définir une dualité parfaite
\[
(11.3.5)\qquad \Phi\otimes\Phi'\longrightarrow\mathbb{Q}/\mathbb{Z}\quad,
\]
ce qui nous donnera en particulier l'accouplement annoncé (11.3.1).
Considérons la suite exacte de $I$-modules
\[
(11.3.6)\qquad 0\longrightarrow U\xrightarrow{i}U\otimes_{\mathbb{Z}_\ell}\mathbb{Q}_\ell
\xrightarrow{j}U\otimes D_\ell\longrightarrow 0\quad,
\]
et la suite exacte de cohomologie associée, en prenant la cohomologie de $I$ définie par les \uline{cochaînes continues} à coefficients dans les groupes topologiques envisagés :
\diagram{
\begin{tikzcd}[column sep=small]
(11.3.7) & 0 \arrow[r] & H^0(I,U) \arrow[r,"i^0"] \arrow[d,phantom,"\wr" description] & H^0(I,U\otimes\mathbb{Q}_\ell) \arrow[r,"j^0"] \arrow[d,phantom,"\wr" description] & H^0(I,U\otimes D_\ell) \arrow[r,"\partial"] \arrow[d,phantom,"\wr" description] & H^1(I,U) \arrow[r,"i^1"] & H^1(I,U\otimes\mathbb{Q}_\ell) \arrow[d,phantom,"\wr" description] \\
& & U^I & (U\otimes\mathbb{Q}_\ell)^I\simeq U^I\otimes_{\mathbb{Z}}\mathbb{Q} & (U\otimes D_\ell)^I & & H^1(I,U)\otimes_{\mathbb{Z}}\mathbb{Q}\ .
\end{tikzcd}
}
Il est immédiat que $\Phi$ défini par (11.3.4) n'est autre que $\operatorname{coker}j^0$, donc est canoniquement isomorphe à $\operatorname{Ker}i^1$, i.e. on trouve une suite exacte canonique
\[
(11.3.8)\qquad 0\longrightarrow\Phi\longrightarrow H^1(I,U)\longrightarrow H^1(I,U)\otimes_{\mathbb{Z}}\mathbb{Q}\quad,
\]
\margin{452}
\margin{452}
\margin{\unsure{141}}
